\section{有界型空间}

在本节中，E 表示一个局部凸空间，B 表示它的典范有界结构（III, 第 3 页, 定义 5）。

引理 1. --- 设 G 是一个半赋范空间，p 是它的半范数，u 是从 G 到 E 中的一个线性映射。下列条件等价：

\begin{enumerate}
\def\labelenumi{(\roman{enumi})}
\item
  \(u\) 连续；
\item
  \(G\) 的单位球在 \(u\) 下的像在 \(\mathbf{E}\) 中是有界的；
\item
  对每个趋于 0 的点序列 \((x_{n})\)（诸点属于 \(\mathbf{G}\)），序列 \((u(x_{n}))\) 在 E 中有界。
\end{enumerate}

立即可见 (i) 蕴含 (ii)（III, 第 4 页, 推论 1），并且 (ii) 蕴含 (iii)。设 V 是 E 中 0 的一个邻域；如果 \(u^{-1}(\mathrm{V})\) 不是 G 中 0 的邻域，那么存在点的一个序列 \((y_n)\)，其诸点属于 \(\mathrm{G} - u^{-1}(\mathrm{V})\)，使得 \(p(y_n) \leqslant \frac{1}{n^2}\)。于是序列 \(x_n = ny_n\) 在 G 中趋于 0，且 \(u(x_n) \notin n\mathrm{V}\)，这蕴含序列 \((u(x_n))\) 不是有界的。因此 (iii) 蕴含 (i)。

命题 1. -- 下列条件等价：

\begin{enumerate}
\def\labelenumi{(\roman{enumi})}
\tightlist
\item
  \(\mathbf{E}\) 上每个在 \(\mathbf{E}\) 的有界子集上有界的半范数都是连续的。
\end{enumerate}

(i') E 的吸收 E 的诸有界子集（I, 第 7 页, 定义 4）的每个凸均衡子集都是 E 中 0 的一个邻域。

\begin{enumerate}
\def\labelenumi{(\roman{enumi})}
\setcounter{enumi}{1}
\tightlist
\item
  E 是诸半赋范空间 \(E_{A}\) 的归纳极限，其中 A 取遍 E 的闭的、凸的、均衡的且有界的子集所成的递增有向集。
\end{enumerate}

(ii') 存在半赋范空间的一个族 \((\mathrm{E}_{i})_{i\in\mathrm{I}}\)，并且对每个 \(i\in I\)，存在一个线性映射 \(u_{i}:E_{i}\to E\)，使得 E 的拓扑是使诸 \(u_{i}\) 连续的最细局部凸拓扑。

\begin{enumerate}
\def\labelenumi{(\roman{enumi})}
\setcounter{enumi}{2}
\tightlist
\item
  对任意的局部凸空间 F，线性映射 \(u: E \to F\) 连续，当且仅当对 E 中趋于 0 的每个点序列 \((x_{n})\)，序列 \((u(x_{n}))\) 在 F 中有界。
\end{enumerate}

(iii') 对任意的半赋范空间 \(\mathbf{F}\)，线性映射 \(u: \mathbf{E} \to \mathbf{F}\) 连续，当且仅当 \(u(\mathbf{X})\) 在 \(\mathbf{F}\) 中有界，其中 \(\mathbf{X}\) 取遍 \(\mathbf{E}\) 中的有界集。

鉴于半范数与凸的、均衡的、吸收的子集之间的对应（II, 第 20 页），立即可见 (i) 与 (i') 等价。若 \(p\) 是 E 上的一个半范数，它在每个 \(\mathrm{E}_{\mathrm{A}}\) 上连续，则 \(p\) 在 E 的有界子集上有界；因而 (i) 蕴含 (ii)（II, 第 27 页, 命题 5）。显然 (ii) 蕴含 (ii')。

现在设 \((\mathrm{E}_{i}, u_{i})_{i\in I}\) 如 (ii') 所述，并设 u 是从 E 到一个局部凸空间 F 中的一个线性映射，使得 \((u(x_{n}))\) 在 F 中有界，对 E 中趋于 0 的每个点序列 \((x_{n})\) 皆然。由 III, 第 11 页 的引理 1 推出，线性映射 \(u\circ u_{i}:E_{i}\to F\) 对所有 \(i\in I\) 连续。因此，若 E 的拓扑是使诸 \(u_{i}\) 连续的最细局部凸拓扑，则 u 连续（II, 第 27 页, 命题 5）。这就证明了 (ii') 蕴含 (iii)。

立即可见 (iii) 蕴含 (iii')（III, 第 3 页, 推论）。最后，若 \(p\) 是 E 上的一个半范数，它在 E 的有界子集上有界，则条件 (iii') 断言恒等映射从 E 到半赋范空间 (E, \(p\)) 中连续；换言之，\(p\) 连续。这就证明了 (iii') 蕴含 (i)。

定义 1. --- 称局部凸空间为有界型空间，如果它满足命题 1 的诸等价条件。

例. --- 1) 每个半赋范空间都是有界型空间。

\begin{enumerate}
\def\labelenumi{\arabic{enumi})}
\setcounter{enumi}{1}
\item
  特别地，每个有限维局部凸空间都是有界型空间。
\item
  由于终局部凸拓扑的传递性（II, 第 28 页, 推论 2），我们由条件 (ii') 立即推出：若 \((\mathrm{E}_i)_{i\in I}\) 是有界型局部凸空间的一个族，且给 E 赋予使诸线性映射 \(u_{i}:\mathrm{E}_{i}\to \mathrm{E}\)（对 \(i\in \mathrm{I}\)）连续的最细局部凸拓扑，则 E 是有界型空间。特别地，有界型空间的归纳极限、直和、商空间都是有界型空间。
\end{enumerate}

另一方面，有界型空间的闭子空间未必是有界型空间（IV, 第 63 页, 习题 8）。

推论. --- 每个 Hausdorff 且半完备的有界型空间都是 Banach 空间的一个归纳极限。

事实上，其中 A 闭且有界的那些空间 \(E_{A}\) 都是 Banach 空间（III, 第 8 页, 推论）。

命题 2. --- 局部凸可度量化空间是有界型空间。

设 E 可度量化，p 是 E 上的一个半范数，它在 E 的有界子集上有界，但不连续。设 A 是所有 \(x \in E\) 中使 \(p(x) < 1\) 者构成的集。设 \((\mathrm{V}_{n})_{n \geqslant 1}\) 是构成 E 中 0 的一个基本邻域系的递减序列。由于 p 不连续，A 不是 0 的邻域；因而对每个 n \textgreater{} 0，有 \(A \not\subset n^{-1}V_{n}\)，且存在一点 \(x_{n}\) 位于 \(V_{n}\) 中，使得 \(n^{-1}x_{n} \notin A\)，即 \(p(x_{n}) \geqslant n\)。序列 \((x_{n})\) 趋于 0，故有界（III, 第 3 页, 推论）；这与关于 p 的假设矛盾。

推论. --- 每个 Fréchet 空间（II, 第 24 页）都是 Banach 空间的一个归纳极限。

